\documentclass[11pt,a4paper]{article}
\usepackage[T1]{fontenc}
\usepackage[utf8]{inputenc}
\usepackage{lmodern,microtype}
\usepackage[margin=.78in]{geometry}
\usepackage{amsmath,amssymb,mathtools,booktabs,tabularx,array,enumitem}
\usepackage[hidelinks]{hyperref}
\usepackage{fancyhdr}
\setlength{\headheight}{14pt}
\pagestyle{fancy}\fancyhf{}
\lhead{\small Sync HotStuff RR, Part III}\rhead{\small Staking / v3}
\cfoot{\thepage}
\setlength{\parindent}{0pt}\setlength{\parskip}{5pt}
\setlength{\emergencystretch}{2em}
\setlist{nosep,leftmargin=*}
\newcolumntype{Y}{>{\raggedright\arraybackslash}X}
\newcommand{\Hash}{\mathsf H}
\newcommand{\tagg}[1]{\mathsf{#1}}
\newcommand{\cfg}{\mathcal C}
\hypersetup{pdftitle={Sync HotStuff RR, Part III: Staking and Authority, v3}}
\begin{document}
\begin{center}
{\Large\bfseries Part III: Staking and Authority}\\[3pt]
{\small Version 3 \quad 19 October 2026}
\end{center}

L1 updates stake in long epochs. Regional sessions can rotate much faster,
using one snapshot many times. This document has three modules: public
selection interfaces, the stake lifecycle, and session authority. Together
they fix the next committee and its initial leader early, while preserving
Part II's fast regular handoff.

Part II derives the append-only global DATA prefix $G$ from finalized
Espresso records and finalized L1 controls. Regions can produce ahead
of $G$. A Force is an authorized branch update: when it takes effect,
old-branch BFT programs stop and a new session starts at the retained
tip of $G$. Accepted DATA is never removed.

To release funds safely, an effective multisig Force closes old BFT
authority and raises a minimum snapshot number. This batches withdrawals
without an L1 transaction at each regular handoff. It can discard the
production suffix beyond global finality; Section~2 states this cost.

\section{Data and selection interfaces}
\label{sec:data}
\subsection{Sessions, blocks and fixed inputs}
A session $s=(\beta,r)$ is one region's production turn. The branch number
$\beta$ changes on Force; $r$ is a turn index, not a region ID. A view is
$(\beta,r,g,t)$, where $g$ counts recoveries and $t$ counts leader terms.
Initially $g=t=0$. For a block $B$, $h(B)$ is its BFT DATA height and
$\lambda(B)$ its authenticated, finalized L1 reference height.

Let $A_s$ be the initial execution anchor and $T_s$ the last DATA of $s$,
at height $h(A_s)+\mathrm{TB}$ for its fixed DATA quota $\mathrm{TB}$.
Part II completes a session on a full terminal FINAL receipt, which
certifies regional commitment. A proposal quorum certificate (QC) alone
does not suffice. A regular successor $s$ of $p$ uses two different blocks:
\[
 \underbrace{E_s=A_p}_{\text{selection reference known when }p\text{ starts}},
 \qquad
 \underbrace{A_s=T_p}_{\text{execution anchor fixed by handoff}}.
\]
The first fixes a plan for synchronization and leader selection. The
second, its receipt, and the exact execution state complete the authority.
Only Part II's Ready check and one-time Enter permit BFT signing.

\begin{tabularx}{\linewidth}{@{}>{\raggedright\arraybackslash}p{.22\linewidth}Y@{}}\toprule
Object & Fixed fields\\\midrule
Stake position $o$ & Permanent ID, region, operator key, withdrawal owner,
escrow $B_o$, request state, exit epoch and processed offenses.\\
Snapshot $\cfg_e$ & Epoch, cutoff, activation height $a_e$, members and
weights, cutoff banned-key set $J_e$, configuration ID and publication proof.\\
Branch policy & Candidate-order rule, boot list, ON/OFF map $S_\beta$,
and minimum snapshot number $m_\beta$.\\
Session authority & Target pair, reference $E_s$, configuration ID, region,
leader rule, execution anchor $A_s$, and authenticated handoff.\\\bottomrule
\end{tabularx}

In $\cfg_e$, $V_i(\cfg_e)$ is region $i$'s member set and $w_o$ a member's
frozen weight. Escrow is separate: after a penalty, weight 100 may be
backed by only 70 remaining tokens; old certificates still use 100.
Genesis fixes quota $\mathrm{TB}>0$, maximum DATA per term $K$, each
region's delay bound $\Delta_i$, and fault fraction $0\le\varphi_i<1/2$.
A session binds these constants and $F_i=\varphi_i\sum_{o\in V_i}w_o$.
Part II assumes actual faulty weight at most $F_i$, with certificate
threshold $>F_i$, throughout that authority's accepted lifetime.

\newpage
\subsection{Snapshot, region and leader functions}
At a verified finalized L1 height $\ell$, define
$e(\ell)=\max\{e:a_e\le\ell\}$. Verify the complete applicable registry
history, not merely the existence of an old snapshot. Missing records or
member data cause a wait. For regular successor $s=(\beta,r+1)$ of
$p=(\beta,r)$, set
\begin{equation}\label{eq:selection}
 E_s=A_p,\qquad h_s^{\rm sel}=h(E_s),\qquad
 e_s=\max\{m_\beta,e(\lambda(E_s))\}.
\end{equation}
The weights come from $\cfg_{e_s}$. The minimum is fixed for the branch.
It changes the anchor-based choice only when necessary to exclude retired
snapshots after a withdrawal Force; it never reads the latest local head.

\paragraph{Region.}
The multisig fixes $\tagg{randList}(\beta,r)$, a deterministic permutation
of registered regions for $p$'s regular successor. Select its first region
$i$ with $S_\beta(i)=\tagg{ON}$ and total snapshot weight $W_i>0$.
If none qualifies, the boundary is pending. Missing evidence is not zero
weight, and silence does not authorize skipping a selected region.
The list rule specifies whether the current region comes last; a separate
boot list selects the first region after genesis or Force. These lists
are not stake-weighted region lotteries.

\paragraph{Weighted choice.}
$\tagg{WeightedRand}(\mathcal W,z)$ is a pure function. Sort positive
integer weights by position ID and let $W$ be their sum. Specified hash
expansion and rejection sampling derive an integer $x\in[0,W)$; choose
the member whose cumulative interval contains $x$. For weights
$(60,30,10)$, $x=73$ selects the second member. An empty set returns $\bot$.
Encoding, hash algorithm, expansion and ID order are protocol constants.

\paragraph{Leader.}
Let $J_{s,g}$ be the fixed banned-key set for generation $g$: $J_{e_s}$
for $g=0$, and the finalized recovery record of Section~3 for $g>0$.
Let $\mathcal W_s^g$ contain the selected region's members not in this set,
with their original weights. The actual leader is
\begin{equation}\label{eq:leader}
 L_s(g,t)=\tagg{WeightedRand}\bigl(\mathcal W_s^g,
 \Hash(\texttt{LEADER},\mathit{chainId},h_s^{\rm sel},\beta,r_s,g,t)\bigr).
\end{equation}
The list uses the \emph{current} turn $r$; the hash uses the \emph{target}
turn $r_s=r+1$. Different views may select the same leader. All voting
members and certificate weights stay fixed, including banned proposers.
The initial schedule is known in advance; a later recovery record may
exclude additional keys, without obtaining a new random seed.

\paragraph{Example.}
In $(0,0)\to(0,1)\to(0,2)$,
$E_{(0,2)}=A_{(0,1)}=T_{(0,0)}$. Suppose this block has DATA height 100
and L1 reference 2500. The seed is 100 and the configuration is
$\cfg_{\max(m_0,e(2500))}$. If $(0,1)$ ends at 120, $(0,2)$ executes
after 120, but keeps the configuration selected from 100. Even if 120
references a newer snapshot, only the plan for $(0,3)$ reads that reference.

Uniform sampling gives probability $w_o/W$ among eligible proposers.
Public height hashing alone does not establish unpredictability,
resistance to stake/identity selection, or schedule fairness. The protocol
uses a reproducible weighted schedule, not a claimed unbiased beacon.

\newpage
\section{Stake lifecycle and liability}
\label{sec:stake}
\subsection{Periodic snapshots and queued changes}
Use one token and whole-position joins and exits. One key belongs to at
most one live position; IDs are never reused. Key or region changes need
new registration after old liability is settled. Genesis fixes the
minimum bond, FIFO stake/count budgets, and entry/exit delays in L1 blocks.
Each admitted position must fit a fresh budget.

Choose a long epoch length $M$ and notice interval $0<P<M$ in L1 blocks.
Snapshot $e$ activates at $a_e=a_0+eM$ with cutoff $c_e=a_e-P$.
At each cutoff, read the complete L1 prefix strictly before $c_e$, apply
the rules below, record the current banned-key set, and freeze $\cfg_e$.
Publish before $a_e$. Late publication cannot change the cutoff; consumers
wait if the required snapshot is unavailable. Genesis supplies $\cfg_0$.

\begin{enumerate}[label=\textbf{\arabic*.},itemsep=3pt]
\item \textbf{Join.} Deposit at least the minimum bond within the admission
cap, register region and keys, prove key possession and owner authorization,
and enter the regional inbound FIFO in L1 order.
The owner may cancel and reclaim a pending join only before any frozen
snapshot includes it; cancellation removes its live queue entry.
\item \textbf{Request exit.} Only a position already included in a frozen
snapshot, with no exit record or pending exit, may enqueue an owner-authorized
irrevocable exit. It remains eligible until processed. A request grants no
permission to stop an assigned duty or withdraw funds.
\item \textbf{Build the next snapshot.} First remove incumbents below
the minimum remaining bond and clear their queue entries. Then select
each remaining FIFO's longest mature prefix
fitting its amount and count budgets. Charge joins their locked balances,
and voluntary exits their weights in $\cfg_{e-1}$. Include surviving
positions with their remaining balances as integer weights.
\item \textbf{Record retirement.} Every omitted incumbent, voluntary or
mandatory, receives its first exclusion epoch $e_x=e$ and an exit record.
It stays absent from every later snapshot. Consume its pending exit
request. Penalty losses and mandatory removals are separate from the
voluntary exit budget. A positive residual balance remains withdrawable
under the same boundary rule below.
\end{enumerate}

Count inflow and outflow separately, even when their net change is zero.
No regional stake floor vetoes exit. Snapshot activation changes future
selection only; it does not edit a fixed session or plan. Queue progress
requires continued L1 processing and positive budgets.

\subsection{Objective slashing}
Evidence binds the chain, protocol, historical position and configuration,
region, session and initial execution anchor. Here $j$ is the fresh block
position within a term. A DRAIN is Part II's recovery block carrying no
new DATA. For one authorized key, the following signed pairs are offenses:

\begin{tabularx}{\linewidth}{@{}>{\raggedright\arraybackslash}p{.24\linewidth}Y@{}}\toprule
Message & Conflict within the same session\\\midrule
PROPOSE & Different admissible DATA/DRAIN proposals at the same $(g,t,j)$,
with Part II admission and quota evidence.\\
Fresh vote & Different block choices at the same $(g,t,j)$; DATA and DRAIN
share the one-vote rule.\\
NEWVIEW / ANCHOR vote & Different recovery starting blocks (openings)
in one generation, even if one extends the other.\\
FINAL & Different DATA hashes at one height; FINAL has no $g,t$.\\\bottomrule
\end{tabularx}

Timeouts, BLAME, OFF status, forwarding evidence, and different QCs for the
same opening are not these offenses. A new branch may legitimately reuse
a DATA height. Section~3 specifies how penalties affect future leaders.

\newpage
\subsection{Penalty execution and a withdrawal boundary}
Anyone may submit a complete offense proof to L1. Verify its signatures,
historical authority and conflict predicate. A valid proof atomically
records the offense, debits a public fraction of the historical weight
capped by remaining escrow, and permanently bans the key in L1's proposer
registry. Invalid or incomplete proofs do not freeze funds. Deduplicate
by position, offense type and conflict slot, not evidence serialization.
Total debits cannot exceed escrow. No discretionary tribunal is needed.

\paragraph{Why exit needs a boundary.}
An exiting node may still belong to a running session or a prepared
successor. Some such sessions are ahead of $G$, so inspecting $G$ does
not enumerate them. This version releases exits in batches through an
existing multisig Force, whose branch-wide closure covers them all.

Extend each branch update with its order policy and requested minimum
snapshot $m$. Along the valid L1 control log, minima never decrease; the
requested snapshot is already activated at the update's publication
height. A matching update's Espresso activation closes the old branch
and installs its new policy, status map and minimum. A stale update
changes none of them. Future branches inherit or raise the minimum.

A position absent from every snapshot starting at $e_x$ gets no new
authority after an effective Force installs $m\ge e_x$: equation
(\ref{eq:selection}) also governs later Force bootstraps. That Force closes
all old BFT programs and cancels their plans, including sessions ahead of
$G$. It does not erase historical penalties or data-service obligations.

\paragraph{The release check.}
Submit the exit record and proof of such an effective Force. From an
authenticated checkpoint, verify complete applicable L1 and Espresso
history, the activation position, and that the target branch was current.
Use full replay of Part II with this module's rules, or a sound proof of
that execution; Espresso inclusion of a claim is insufficient. A reusable
checkpoint binds the complete replay state and source commitments:
$G$; the NEXT/PENDING tail (the next DATA owner or unassigned boundary);
status, order policy and minimum; and the position reached in the ordered
control log with its pending updates, at the exact Espresso position.

Record the first verified clearance for the position once. Repeated
proofs neither backdate nor restart its clock. Let $t_c$ be the real time
when L1 accepts it. Part II bounds closure delivery by $T_E$, then delivery
of available complete results by $T_{\rm pub}+T_E$. Let $T_{\rm ev}$ cover
the additional investigation, collection and submission of relevant
evidence, including discarded sessions. Assume valid slash transactions
finalize within $T_{\rm L1}^{\rm stake}$ of actual submission. Choose
\[
 T_{\rm hold}\ge 2T_E+T_{\rm pub}+T_{\rm ev}
                    +T_{\rm L1}^{\rm stake}+\varepsilon_{\rm clock}.
\]
The first $T_E$ allows nodes to learn closure; $T_{\rm pub}+T_E$ covers
their remaining result delivery. The deployed L1 timelock and clock
margin must prevent expiry before this real interval. These are explicit
service/clock assumptions, not consequences of an average block time.

Release residual escrow once, strictly after the deadline and after
applying valid proofs accepted by it. A finite window cannot catch evidence
hidden forever or unbounded offline investigation. Data must be retained
or durably transferred under Part II before the operator departs.

\newpage
\section{Deriving and using authority}
\label{sec:authority}
\subsection{Initial setup and regular handoff}
Genesis fixes the chain and algorithm identifiers, registered regions,
the BFT and staking parameters of Sections~1--2, the initial policy and
$S_0$, $m_0=0$, and an authenticated DATA anchor $A_0$ with
$\lambda(A_0)\ge a_0$. For $(0,0)$ set $E=A=A_0$ and use the boot list
and snapshot rule. The checkpoint authenticates the complete initial
L1/Espresso replay state. This is the base case, not a regular predecessor.

For each running $p=(\beta,r)$:
\begin{enumerate}[label=\textbf{\arabic*.},itemsep=4pt]
\item \textbf{Prepare.} Use $E_s=A_p$ to select the next snapshot, region
and initial leader. Fix a plan for $s=(\beta,r+1)$. Its committee follows
$p$ and prepares state, but issues no BFT signatures for $s$.
\item \textbf{Bind the handoff.} On a complete terminal receipt, fix
$A_s=T_p$ and verify that the branch and plan remain applicable. The initial
$\tagg{AUTH}_{\beta,r_s,0}$ binds the plan, both blocks, the snapshot and
policy IDs, the handoff, and the genesis-derived BFT parameters.
\item \textbf{Enter.} Validate the exact state after $A_s$ and all sources.
Part II's Ready check and guarded, one-time Enter start $s$. Its first
leader directly proposes the child of $A_s$. A later L1 snapshot does
not replace the plan. No new L1 transaction or global-finality wait is
added to this handoff.
\end{enumerate}

The session keeps its voting snapshot, region, selection reference, quota
and seed across all terms and recoveries. Its execution opening may change
during recovery, but that new opening never chooses the recovery leader.

\subsection{Recovery authority and proposer exclusion}
A quit proof is Part II's certificate of insufficient progress or valid
leader-equivocation evidence. For $g>0$, anyone can submit the complete
quit proof for $(\beta,r,g-1)$ and its authority dependencies to L1.
The contract creates at most one immutable
\[
 \tagg{REC}(\beta,r,g)=(\mathit{authorityId}_s,g,J_{s,g}),
\]
with the banned-key set exactly as it stands before that transaction.
The quit proof must justify entry to $g$; a bare counter cannot reserve
future records. Wait for record finality before entering the generation,
as Part II already waits for recovery authority. The initial session
context and this record determine $\tagg{AUTH}_{\beta,r,g}$; then use
(\ref{eq:leader}) on the eligible proposers. If none remains, wait for
external recovery; do not change the voting quorum.

Initial generation uses the selected snapshot's cutoff set $J_{e_s}$.
A later penalty cannot rewrite its frozen initial schedule. Later recovery
records and newly cut snapshots apply the permanent ban. Exclusion progress
requires proof delivery and the corresponding authority updates; public
hash schedule fairness remains a separate Part II service premise.

\subsection{Force bootstrap}
At an effective Force, Part II fixes the retained tip of $G$ and target
$r$, including its zero-output reuse rule. For the first new session, use
that tip as both $E_s$ and $A_s$, the new boot list and minimum, and the
same snapshot/leader functions. Subsequent regular handoffs again use the
predecessor's initial anchor. Old plans, votes, locks and timers do not
carry over, even when the turn number is reused. The new session uses the
verified execution state at its retained anchor.

\newpage
\subsection{Part II interfaces and the resulting guarantees}
Part II v4.6 indexes authority contexts by $(\beta,r,g)$ and derives the
leader for the full view $(\beta,r,g,t)$. Here generation zero is the
initial session context; recovery combines it with $\tagg{REC}(\beta,r,g)$.
All generations keep the same configuration ID and voting snapshot.
No new L1 transaction is required for a regular handoff or ordinary term
change. FINAL remains session-bound, without $g$ or $t$.
Candidate and boot lists are this module's replacement for Part II's
fixed cyclic NextLive selector. Admission, handoff verification and
public replay use the same authority construction; quit, lock and timer
rules are unchanged.

For first DATA, the L1 reference lower bound is the maximum of the
incoming anchor's reference, snapshot publication and activation heights,
and the applicable branch-control publication height. Genesis supplies
its initial control bound. This permits an old retained execution anchor
while ensuring fresh proposals can cite the authority that authorizes them.

\paragraph{Deterministic authorization.}
Starting from genesis, authenticated history fixes each reference and
branch policy. The snapshot rule gives one configuration; ordered
eligibility gives one region or a pending boundary. The initial leader
uses its snapshot's fixed banned-key set. Each later generation uses its
unique finalized recovery record. Thus every complete, admitted view has
the same leader at every verifier. Missing dependencies cause a wait,
not a competing assignment. A plan alone never permits entry.

\paragraph{Voluntary withdrawal preserves the stated liability.}
Before clearance, escrow remains locked. A qualifying effective Force
closes every older branch, including all its speculative programs and
plans. Its minimum, inherited by future branches, excludes the position
from every new authority. The holding period then covers closure delivery,
remaining result delivery and timely evidence processing under the stated
bounds. Historical certificates retain frozen weights; prior penalties
are not counted again as remaining collateral. This is not a promise to
pay for every ever-hidden offense or to replenish slashed funds.

\paragraph{External history.}
Part II obtains $G$ by replaying finalized Espresso inputs and complete
L1 controls. Raw Espresso inclusion is not acceptance into $G$; once
accepted, DATA cannot be removed by Force. Withdrawal relies on that
external history and sound verification, not a new vote by exited regional
keys. No extra social checkpoint mechanism is introduced for this layer.
Existing L1/Espresso trust anchors and historical data availability remain
assumptions. An old key cannot reopen a closed, retired authority.

\paragraph{Accepted costs and limits.}
\begin{itemize}[itemsep=3pt]
\item The initial schedule is public and predictable. Its stake-weighted
construction does not prove resistance to identity selection or targeting.
\item Adoption can lag during a long session. Final withdrawal needs a
multisig batch Force and external services; it has no unconditional time
bound. The Force may discard production beyond $G$.
\item Recovery still waits for L1. Stake losses do not reweight a running
committee. Each region's frozen-weight fault bound and sufficient
uninterrupted operation remain Part II assumptions.
\item Whole-position exits keep the first version small. Partial exits,
delegation and changes to the fixed protocol parameters need explicit
extensions rather than silently changing these rules.
\end{itemize}

\paragraph{References.}
Part II v4.6: session/generation authority contexts, full-view leaders,
FINAL, Force-close and service bounds.
Part I: recovery authority, proposer exclusion and conditional progress.
For queue/lifecycle patterns, see Ethereum's
\href{https://ethereum.github.io/consensus-specs/specs/fulu/beacon-chain/}{Fulu pending deposits}
and \href{https://ethereum.github.io/consensus-specs/specs/electra/beacon-chain/}{Electra exit and slashing rules}.
Their withdrawal constants are not imported here.
\end{document}
